\documentclass{article}
\usepackage{amsmath, amssymb}

\title{Periodicity}
\author{}
\date{}

\begin{document}
	
	\maketitle
	
	\begin{enumerate}
		
		\item (P1) (O1) Let $P \in \mathbb{Z}[x]$ and $a_1 = 1$, with $a_{n+1} = P(a_n)$. Define $b_n \equiv a_n \pmod{N}$. Show that $b_n$ is eventually periodic.
		
		\item (P2) (O2) Let $a_1 = 1$ and $a_{n+1} = a_n^2 + a_n - 1$. Show that
		\[
		\gcd(n, a_n) = 1
		\]
		for all $n$.
		
		\item (P3) (O1) Let $p(n)$ be the greatest prime divisor of $n$. Define the sequence by $a_1 \in \mathbb{N}$ and
		\[
		a_{n+1} = p(a_n) + 1.
		\]
		Show that $a_n$ is eventually periodic and determine the length of the period.
		
		\item (P4) (O2) Let
		\[
		\operatorname{rad}(n) = \prod_{p \mid n} p.
		\]
		Show that the sequence defined by $a_1 \in \mathbb{N}$ and
		\[
		a_{n+1} = \operatorname{rad}(a_n) + 1
		\]
		is eventually periodic and determine the length of the period.
		
		\item (P5) (O3) Let $P \in \mathbb{Z}[x]$ and $a_1 \in \mathbb{Z}$, with
		\[
		a_{n+1} = P(a_n).
		\]
		Suppose that the sequence $(a_n)$ is periodic. Show that the length of the period $L$ satisfies $L \le 2$.
		
	\end{enumerate}
	
\end{document}